%------------------------------------------------------------------------------%
\documentclass[fleqn,aspectratio=169,12pt]{beamer}

\usepackage{gaal}

\title{Hipérbole}

%------------------------------------------------------------------------------%
\begin{document}

\addtitle

%------------------------------------------------------------------------------%
\section{Equação Canônica da Hipérbole}

%------------------------------------------------------------------------------%
\begin{frame}{Equação Canônica da Hipérbole}
  Equação \emph{canônica da hipérbole}
  \[
    \frac{x^2}{a^2} - \frac{y^2}{b^2} = 1
  \]
  com \(a > 0\) e \(b > 0\)

  \bigskip
  Hipérbole centrada na origem e alinhada com os eixos $x$ e $y$
\end{frame}

%------------------------------------------------------------------------------%
\begin{frame}{Equação Matricial da Hipérbole}
  Equação \emph{matricial da hipérbole}
  \[
    v^t A v = 1
  \]
  com
  \[
    v = \begin{bmatrix}
      x \\
      y
    \end{bmatrix}
    \qquad\qquad
    A = \begin{bmatrix}
      \dfrac{1}{a^2} & 0              \\[2mm]
      0              & -\dfrac{1}{b^2}
    \end{bmatrix}
    \qquad\qquad
    a > 0
    \qquad
    b > 0
  \]
\end{frame}

%------------------------------------------------------------------------------%
\begin{proofframe}{Pontos Relevantes da Hipérbole}
\begin{columns}[t]
\column{0.4\textwidth}
  Fazendo \(y=0\)
  \begin{align*}
    \frac{x^2}{a^2} - \frac{y^2}{b^2} & = 1     \\
    \frac{x^2}{a^2} - \frac{0}{b^2}   & = 1     \\
    \frac{x^2}{a^2}                   & = 1     \\
    x^2                               & = a^2   \\
    \abs{x}                           & = a     \\
    x                                 & = \pm a
  \end{align*}

\column{0.4\textwidth}
  Fazendo \(x=0\)
  \begin{align*}
    \frac{x^2}{a^2} - \frac{y^2}{b^2} & = 1     \\
    \frac{0}{a^2} - \frac{y^2}{b^2}   & = 1     \\
    -\frac{y^2}{b^2}                  & = 1     \\
    y^2                               & = -b^2  \\
    \nexists\, y \in \R
  \end{align*}
\end{columns}
\end{proofframe}

%------------------------------------------------------------------------------%
\begin{frame}{Elementos Geométricos da Hipérbole}
  \emph{Centro} da hipérbole
  \[
    C = (0, 0)
  \]

  Relação fundamental da hipérbole: os focos são
  \[
    F_1 = (-c, 0) \qquad
    F_2 = ( c, 0)
    \qquad\qquad
    c^2 = a^2 + b^2
  \]

  A hipérbole não cruza o eixo $y$ e cruza o eixo $x$ nos seus vértices
  \[
    (-a, 0)
    \qquad
    (a, 0)
  \]

  A hipérbole é simétrica em relação aos eixos

  A hipérbole possui assintotas
\end{frame}

%------------------------------------------------------------------------------%
\begin{frame}{Gráfico da Hipérbole}{}
{\color{red} Incluir gráficos}
\end{frame}

%------------------------------------------------------------------------------%
%------------------------------------------------------------------------------%
\begin{question}
  Dada a equação da hipérbole
  \[
    \frac{x^2}{9} - \frac{y^2}{4} = 1
  \]
  \vspace{-1.5\baselineskip}
  \begin{enumerate}
    \item encontre os vértices
    \item escreva a equação matricial
    \item esboce a curva
  \end{enumerate}
\end{question}

%------------------------------------------------------------------------------%
\begin{resolution}{Vértices}
  Comparando
  \[
    \frac{x^2}{9} - \frac{y^2}{4} = 1
  \]
  \pause
  com a equação canônica
  \[
    \frac{x^2}{a^2} - \frac{y^2}{b^2} = 1
  \]
  \pause
  identificamos
  \[
    \pause
    a^2 = 9
    \qquad
    \pause
    a = 3
    \qquad\qquad\qquad
    \pause
    b^2 = 4
    \qquad
    \pause
    b = 2
  \]
  \pause
  os vértices são os pontos
  \[
    \pause
    (-a, 0) = (-3, 0)
    \qquad\qquad
    \pause
    (a, 0) = (3, 0)
  \]
\end{resolution}

%------------------------------------------------------------------------------%
\begin{resolution}{Equação Matricial}
  Equação Matricial
  \[
    v^t A v = 1
  \]
  \pause
  com
  \[
    \pause
    A
    \pause
    = \begin{bmatrix}
      \dfrac{1}{a^2} & 0              \\[2mm]
      0              & -\dfrac{1}{b^2}
    \end{bmatrix}
    \pause
    = \begin{bmatrix}
      \dfrac{1}{9} & 0            \\[2mm]
      0            & - \dfrac{1}{4}
    \end{bmatrix}
  \]
\end{resolution}

%------------------------------------------------------------------------------%
\begin{resolution}{Esboço}
  \centering
  \includegraphics{I-06-Hiperbole-ex-1}
\end{resolution}

%------------------------------------------------------------------------------%
\endinput


%------------------------------------------------------------------------------%
\section{Translação da Hipérbole}

%------------------------------------------------------------------------------%
\begin{frame}{Translação da Hipérbole}
  Novo centro \(C = (h, k)\)

  Novas coordenadas
  \[
    X = x - h
    \qquad
    Y = y - k
  \]

  Equação canônica nas novas coordenadas
  \[
    \frac{X^2}{a^2} - \frac{Y^2}{b^2} = 1
  \]
  \[
    \frac{(x-h)^2}{a^2} - \frac{(y-k)^2}{b^2} = 1
  \]
\end{frame}

%------------------------------------------------------------------------------%
\begin{proofframe}{Equação Matricial da Elipse Transladada}
  Novo de centro
  \[
    C
    = \begin{bmatrix}
      h \\
      k
    \end{bmatrix}
  \]

  \bigskip
  Novas coordenadas
  \[
    V
    = \begin{bmatrix}
      X \\
      Y
    \end{bmatrix}
    = v - C
    = \begin{bmatrix}
      x \\
      y
    \end{bmatrix}
    -
    \begin{bmatrix}
      h \\
      k
    \end{bmatrix}
    = \begin{bmatrix}
      x - h \\
      y - k
    \end{bmatrix}
  \]
\end{proofframe}

%------------------------------------------------------------------------------%
\begin{proofframe}{Equação Matricial da Hipérbole Transladada}
  \begin{align*}
    \emph{V^t A V}    & = \emph{1} \\
    (v-C)^t A (v-C)   & = 1 \\
    \begin{bmatrix}
      x - h \\
      y - k
    \end{bmatrix}^t
    \begin{bmatrix}
      \sfrac{1}{a^2} & 0              \\
      0              & \sfrac{-1}{b^2}
    \end{bmatrix}
    \begin{bmatrix}
      x - h \\
      y - k
    \end{bmatrix} & = 1 \\
    \emph{\frac{(x-h)^2}{a^2} - \frac{(y-k)^2}{b^2}} & = \emph{1} \\
    \frac{b^2(x-h)^2 - a^2(y-k)^2}{a^2b^2}    & = 1 \\
    b^2(x-h)^2 - a^2(y-k)^2                   & = a^2b^2
  \end{align*}
\end{proofframe}

%------------------------------------------------------------------------------%
\begin{proofframe}{Equação Matricial da Hipérbole Transladada}
  \begin{align*}
    b^2(x-h)^2 - a^2(y-k)^2                             & = a^2b^2 \\
    b^2(x^2 -2hx + h^2) - a^2(y^2 - 2ky + k^2)          & = a^2b^2 \\
    b^2x^2 - 2b^2hx + b^2h^2 - a^2y^2 + 2a^2ky - a^2k^2 & = a^2b^2 \\
    b^2x^2 - a^2y^2 - 2b^2hx + 2a^2ky + b^2h^2 - a^2k^2 & = a^2b^2 \\
    \begin{bmatrix}
      x \\
      y
    \end{bmatrix}^t
    {\color{magenta}\begin{bmatrix}[c]
      b^2 & 0   \\
      0   & -a^2
    \end{bmatrix}}
    \begin{bmatrix}
      x \\
      y
    \end{bmatrix}
    +
    {\color{purple}\begin{bmatrix}
      - 2b^2h \\
        2a^2k
    \end{bmatrix}^t}
    \begin{bmatrix}
      x \\
      y
    \end{bmatrix}
    + {\color{violet} (b^2h^2 - a^2k^2 - a^2b^2)} & = 0 \\
    v^t {\color{magenta} M} v
    + {\color{purple} B^t}v
    + {\color{violet} f} & = 0
  \end{align*}
\end{proofframe}

%------------------------------------------------------------------------------%
%------------------------------------------------------------------------------%
\begin{question}
  Dada a equação da hipérbole
  \[
    \frac{(x-1)^2}{16} - \frac{(y+2)^2}{9} = 1
  \]
  \vspace{-1.5\baselineskip}
  \begin{enumerate}
    \item encontre o centro e os vértices
    \item escreva a equação matricial
    \item esboce a curva
  \end{enumerate}
\end{question}

%------------------------------------------------------------------------------%
\begin{resolution}{Centro}
  Comparando
  \[
    \frac{(x-1)^2}{16} - \frac{(y+2)^2}{9} = 1
  \]
  com a equação canônica transladada
  \[
    \frac{(x-h)^2}{a^2} - \frac{(y-k)^2}{b^2} = 1
  \]
  obtemos
  \[
    a^2 = 16
    \qquad
    a = 4
  \]
  \[
    b^2 = 9
    \qquad
    b = 3
  \]
  o centro é
  \(
    C
    = (h, k)
    = (1, -2)
  \)
\end{resolution}

%------------------------------------------------------------------------------%
\begin{resolution}{Vértices}
\begin{columns}[t]
\column{0.6\textwidth}
  Novas coordenadas
  \[
    V
    = \begin{bmatrix}
      x - h  \\
      y - k
    \end{bmatrix}
    = \begin{bmatrix}[c]
      x - 2  \\
      y - (-1)
    \end{bmatrix}
    = \begin{bmatrix}
      x - 2  \\
      y + 1
    \end{bmatrix}
  \]
\column{0.4\textwidth}
  Vértices
  \[
    V
    = \begin{bmatrix}
      x - 2  \\
      y + 1
    \end{bmatrix}
    = \begin{bmatrix}
      \pm 4  \\
      0
    \end{bmatrix}
  \]
\end{columns}

\medskip
Resolvendo
\vspace{-1.5\baselineskip}
\begin{columns}[t]
\column{0.2\textwidth}
  \begin{gather*}
    y + 1 = 0 \\
    y     = -1
  \end{gather*}
\column{0.2\textwidth}
  \begin{gather*}
    x - 2 = 4     \\
    x     = 4 + 2 \\
    x     = 6
  \end{gather*}
\column{0.2\textwidth}
  \begin{gather*}
    x - 2 = -4 \\
    x     = -4 + 2 \\
    x     = -2
  \end{gather*}
\end{columns}

  \medskip
  Vértices:
  \(
    \quad
    \mathcal{V}_1 = (-2, -1)
    \quad
    \mathcal{V}_1 = (6, -1)
  \)
\end{resolution}

%------------------------------------------------------------------------------%
\begin{resolution}{Equação Matricial}
  Equação Matricial
  \[
    V^t A V = 1
  \]
  com
  \[
    A = \begin{bmatrix}
      \dfrac{1}{16} & 0            \\[2mm]
      0             & -\dfrac{1}{9}
    \end{bmatrix}
    \qquad\qquad
    V = \begin{bmatrix}
      x - 2  \\
      y + 1
    \end{bmatrix}
  \]
\end{resolution}

%------------------------------------------------------------------------------%
\begin{resolution}{Esboço}
\end{resolution}

%------------------------------------------------------------------------------%
\endinput


%------------------------------------------------------------------------------%
\begin{question}
  Dada a equação da hipérbole
  \[
    \frac{(x-1)^2}{16} - \frac{(y+2)^2}{9} = 1
  \]
  \vspace{-1.5\baselineskip}
  \begin{enumerate}
    \item encontre o centro e os vértices
    \item escreva a equação matricial
    \item esboce a curva
  \end{enumerate}
\end{question}

%------------------------------------------------------------------------------%
\begin{resolution}{Centro}
  Comparando
  \[
    \frac{(x-1)^2}{16} - \frac{(y+2)^2}{9} = 1
  \]
  \pause
  com a equação transladada
  \[
    \frac{(x-h)^2}{a^2} - \frac{(y-k)^2}{b^2} = 1
  \]
  \pause
  identificamos
  \[
    \pause
    a^2 = 16
    \qquad
    \pause
    a = 4
    \qquad\qquad\qquad
    \pause
    b^2 = 9
    \qquad
    \pause
    b = 3
  \]
  \pause
  Centro: \quad
  \(
    \pause
    C
    = (h, k)
    \pause
    = (1, -2)
  \)
\end{resolution}

%------------------------------------------------------------------------------%
\begin{resolution}{Vértices}
\begin{columns}[t]
\column{0.6\textwidth}
  \onslide<+->{Novas coordenadas}
  \[
    \onslide<+->{V
    = \begin{bmatrix}
      x - h  \\
      y - k
    \end{bmatrix}}
    \onslide<+->{= \begin{bmatrix}[c]
      x - 1 \\
      y - (-2)
    \end{bmatrix}}
    \onslide<+->{= \begin{bmatrix}
      x - 1  \\
      y + 2
    \end{bmatrix}}
  \]
\column{0.4\textwidth}
  \onslide<+->{Condição nos vértices}
  \[
    \onslide<+->{V
    = \begin{bmatrix}
      x - 1  \\
      y + 2
    \end{bmatrix}
    = \begin{bmatrix}
      \pm 4  \\
      0
    \end{bmatrix}}
  \]
\end{columns}

\medskip
\onslide<+->{Resolvendo a condição nos vértices}
\vspace{-1.5\baselineskip}
\begin{columns}[t]
\column{0.2\textwidth}
  \begin{gather*}
    \onslide<+->{y + 2 = 0 \\}
    \onslide<+->{y     = - 2 }
  \end{gather*}
\column{0.2\textwidth}
  \begin{gather*}
    \onslide<+->{x - 1 = 4     \\}
    \onslide<+->{x     = 4 + 1 \\}
    \onslide<+->{x     = 5}
  \end{gather*}
\column{0.2\textwidth}
  \begin{gather*}
    \onslide<+->{x - 1 = -4 \\}
    \onslide<+->{x     = -4 + 1 \\}
    \onslide<+->{x     = -3}
  \end{gather*}
\end{columns}

  \medskip
  \onslide<+->{Vértices:}
  \(
    \quad
    \onslide<+->{\mathcal{V}_1 = (-3, -2)}
    \quad
    \onslide<+->{\mathcal{V}_1 = (5, -2)}
  \)
\end{resolution}

%------------------------------------------------------------------------------%
\begin{resolution}{Equação Matricial}
  Equação Matricial
  \[
    \pause
    V^t A V = 1
  \]
  \pause
  com
  \[
    \pause
    A
    = \begin{bmatrix}
      \dfrac{1}{16} & 0            \\[2mm]
      0             & -\dfrac{1}{9}
    \end{bmatrix}
    \qquad\qquad
    \pause
    V = \begin{bmatrix}
      x - 2  \\
      y + 1
    \end{bmatrix}
  \]
\end{resolution}

%------------------------------------------------------------------------------%
\begin{resolution}{Esboço}
  \centering
  \includegraphics{I-06-Hiperbole-ex-3}
\end{resolution}


%------------------------------------------------------------------------------%
\endinput



%%------------------------------------------------------------------------------%
%\section{Lista Mínima}
%
%%------------------------------------------------------------------------------%
%\begin{frame}{Lista Mínima}
%  \begin{enumerate}
%    \item
%  \end{enumerate}
%
%  \vfill
%  Atenção: A prova é baseada no livro, não nas apresentações
%\end{frame}

%------------------------------------------------------------------------------%
\end{document}
%------------------------------------------------------------------------------%
